\section{Related Work}
\label{sec:related-work}

\paragraph{POMDP and belief-space planning.}
POMDP methods provide the standard formalism for acting under hidden state \citep{Kaelbling1998, Spaan2012, Lauri2023}, and online planners such as POMCP \citep{Silver2010} demonstrate the value of maintaining beliefs over latent state. In manipulation, belief-space replanning can also be necessary for long-horizon tasks with uncertainty \citep{Garrett2020}. This literature establishes that hidden state should be reasoned about explicitly, but it usually targets full inference or planning over latent state rather than a systems-level triage question: before escalating to belief estimation, can a \emph{memory design diagnostic layer} first determine whether the failure is merely hidden previously visible state that only needs to be persisted?

\paragraph{Learned memory architectures.}
Another line of work equips policies with implicit learned memory through recurrent networks \citep{Hausknecht2015}, state-space models \citep{Gu2022}, transformers \citep{Vaswani2017, Brohan2023}, sequence modeling \citep{Chen2021}, world models \citep{Ha2018, Hafner2020}, image-based RL \citep{Yarats2021}, and meta-learning under partial observability \citep{Zintgraf2020}. These methods are often powerful because they combine representation learning, temporal abstraction, and memory in one end-to-end policy. However, that strength also makes diagnosis harder: when performance improves, it is often unclear whether the gain comes from genuinely retaining missing information, from larger model capacity, or from broader changes in training dynamics. What remains underspecified is the diagnostic layer that would justify invoking these heavy mechanisms in the first place, rather than first testing whether explicit low-cost persistence already removes the failure mode.

\paragraph{Explicit cache and goal memory.}
Work on explicit history modeling shows that structured access to past observations can improve manipulation \citep{Guhur2022}, and RoboMME suggests that different memory formats help on different task families \citep{Dai2026}. A related thread studies explicit goal caching or latent goal persistence \citep{Yarats2022, Lynch2020}, which is especially relevant when the missing variable was visible earlier and is compact enough to store directly. This line is closest to our setting because it points to a lightweight memory mechanism that is interpretable and controllable. Yet prior work typically evaluates such memory as one design option among many, rather than positioning it inside a diagnostic layer that asks whether a flat cache should be the first intervention and whether richer memory organization is warranted only after that cache fails.

\paragraph{Recovery policies.}
Recovery policies and retry mechanisms can improve robustness in sequential manipulation under noise and partial observability \citep{Vats2024, Vats2023}. They provide a different axis of adaptation: instead of preserving hidden state better, they alter control behavior after failure or uncertainty is detected. Such methods can be highly effective, but they also consume time, actions, and finite-horizon budget, so a stronger recovery policy is not automatically a lower-cost or cleaner fix than better state persistence. This creates another missing piece for a memory design diagnostic layer: before adding recovery logic, one should know whether the underlying failure is actually a persistence problem that a simple cache would solve more directly.

\paragraph{Gap.}
Taken together, prior work offers strong tools for hidden-state inference, learned memory, explicit caching, and recovery, but not a single \emph{memory design diagnostic layer} that orders them by necessity. The missing question is not simply ``which memory architecture scores highest?'' but ``what is the smallest intervention that plausibly addresses this failure mode?'' Our paper fills that gap by inserting an explicit diagnostic layer ahead of architecture escalation: first test whether hidden previously visible state is the relevant bottleneck, then test whether flat goal persistence already closes the gap, then ask whether richer organization or recovery adds anything beyond that cache, and finally mark the boundary where the missing state has moved beyond goal-only memory. This decision-oriented framing is the main distinction from prior memory or recovery comparisons.
